%! TEX root = main.tex

\section{Studies on Boundary Conditions and Friction}
\label{sc:25}

Two optional modifications of the periodic shallow-water model are examined: the replacement of periodicity by reflecting wall boundary conditions and the introduction of a linear friction term in the momentum equation. In both cases, the theoretical implications for the equivalent second-order equation, the energy balance and the discrete block formulation are derived first. The corresponding numerical evidence is then collected and interpreted in a single dedicated subsection.

\subsection{Part A: Reflecting Wall Boundary Conditions}
\label{ssc:251}

\paragraph{Strong formulation and equivalent second-order condition}
The periodic boundary conditions are replaced by rigid-wall conditions, namely

\begin{equation}
q(0,t)=0, \qquad q(L,t)=0
\label{eq:reflecting_bc_hw2}
\end{equation}

These impose zero discharge at both endpoints, describing a closed basin where no flux enters or leaves the domain. The equivalent condition for the second-order wave equation is obtained by evaluating the momentum equation \(q_t + gH\eta_x = 0\) at \(x=0\) and \(x=L\). Since \(q(\cdot,t)\) vanishes at both walls for every \(t\), its time derivative also vanishes there. Therefore:

\begin{equation}
q_t(0,t)=0 \quad\Longrightarrow\quad gH\eta_x(0,t)=0 \quad\Longrightarrow\quad \eta_x(0,t)=0
\label{eq:neumann_left_wall}
\end{equation}

and, analogously,

\begin{equation}
q_t(L,t)=0 \quad\Longrightarrow\quad gH\eta_x(L,t)=0 \quad\Longrightarrow\quad \eta_x(L,t)=0
\label{eq:neumann_right_wall}
\end{equation}

Hence, homogeneous Dirichlet conditions on \(q\) are equivalent to homogeneous Neumann conditions on \(\eta\). Physically, the wave is reflected by the wall without changing sign.

\paragraph{Effect on the energy balance}
The conserved energy introduced in section \ref{sc:22} is

\begin{equation}
E(t) = \int_0^L \left(\frac{1}{2}g\eta^2 + \frac{1}{2H}q^2\right) dx
\label{eq:energy_rw}
\end{equation}

Repeating the computation carried out for the periodic case yields

\begin{equation}
\frac{dE}{dt}
= -g\left[\eta(x,t) q(x,t)\right]_{x=0}^{x=L}
= -g\Bigl(\eta(L,t)q(L,t)-\eta(0,t)q(0,t)\Bigr)
\label{eq:energy_wall_boundary_term}
\end{equation}

Substituting \eqref{eq:reflecting_bc_hw2}, the boundary contribution vanishes identically:

\begin{equation}
\frac{dE}{dt}=0
\label{eq:energy_reflecting_conserved}
\end{equation}

Reflecting walls therefore preserve the total mechanical energy exactly. Unlike the periodic setting, energy is no longer transported across the boundary; instead, it remains trapped inside the domain and is redistributed through successive reflections.

\paragraph{Effect on the weak formulation and discrete spaces}
The function spaces must be adapted. The periodicity constraint is removed, so the free-surface elevation is sought in

\begin{equation}
V = H^1(0,L)
\label{eq:wall_space_eta}
\end{equation}

Since the discharge must satisfy \eqref{eq:reflecting_bc_hw2} strongly, the corresponding trial and test space becomes

\begin{equation}
V_0 = H_0^1(0,L) = \left\{v\in H^1(0,L): v(0)=v(L)=0\right\}
\label{eq:wall_space_q}
\end{equation}

The weak problem reads: find \(\eta(\cdot,t)\in V\) and \(q(\cdot,t)\in V_0\) such that, for all \(v\in V\) and \(w\in V_0\),

\begin{equation}
\int_0^L \eta_t v  dx + \int_0^L q_x v  dx = 0
\label{eq:wall_weak_1}
\end{equation}
\begin{equation}
\int_0^L q_t w  dx + gH\int_0^L \eta_x w  dx = 0
\label{eq:wall_weak_2}
\end{equation}

The algebraic block structure remains the same as in section \ref{sc:24}; only the treatment of the boundary degrees of freedom changes. In practice, the periodic identification of the first and last node is removed and the endpoint degrees of freedom associated with \(q\) are constrained to zero.

\subsection{Part B: Linear Friction}
\label{ssc:252}

\paragraph{Modified strong formulation}
A linear friction term is added to the momentum equation with damping coefficient \(\gamma > 0\):

\begin{equation}
\begin{cases}
\eta_t + q_x = 0\\
q_t + gH\eta_x = -\gamma q
\end{cases}
\qquad \gamma>0
\label{eq:friction_system_hw2}
\end{equation}

The term \(-\gamma q\) models a resistive effect that opposes the discharge and progressively removes mechanical energy from the system.

\paragraph{Equivalent damped wave equation}
The continuity equation is differentiated with respect to time:

\begin{equation}
\eta_{tt} = -q_{xt}
\label{eq:eta_tt_from_continuity}
\end{equation}

The momentum equation is differentiated with respect to \(x\):

\begin{equation}
q_{tx} + gH\eta_{xx} = -\gamma q_x
\label{eq:momentum_x_diff}
\end{equation}

By Schwarz's theorem, \(q_{xt}=q_{tx}\). Using \eqref{eq:eta_tt_from_continuity} and the identity \(q_x=-\eta_t\) from the continuity equation:

\begin{equation}
\eta_{tt} + \gamma\eta_t = gH\eta_{xx}
\label{eq:telegrapher_hw2}
\end{equation}

This is the Telegrapher's equation. The additional first-order time derivative introduces amplitude attenuation while leaving the wave propagation mechanism otherwise unchanged.

\paragraph{Effect on the energy identity}
Starting from \eqref{eq:energy_rw}, the time derivative of the energy is

\begin{equation}
\frac{dE}{dt}
= \int_0^L \left(g\eta\eta_t + \frac{1}{H}q q_t\right) dx
\label{eq:energy_time_derivative_friction}
\end{equation}

Substituting the modified system \eqref{eq:friction_system_hw2}:

\begin{equation}
\frac{dE}{dt}
= \int_0^L \left[g\eta(-q_x) + \frac{1}{H}q(-gH\eta_x-\gamma q)\right]dx
\label{eq:energy_substitution_friction}
\end{equation}

Collecting terms and recognizing the total derivative \(\eta q_x + q\eta_x = \frac{d}{dx}(\eta q)\):

\begin{equation}
\frac{dE}{dt}
= -g\left[\eta q\right]_0^L - \frac{\gamma}{H}\int_0^L q^2  dx
\label{eq:energy_friction_boundary_plus_bulk}
\end{equation}

Under either periodic or reflecting-wall boundary conditions, the boundary term vanishes, yielding:

\begin{equation}
\frac{dE}{dt} = -\frac{\gamma}{H}\int_0^L q^2  dx \leq 0
\label{eq:energy_dissipation_hw2}
\end{equation}

Friction therefore destroys the conservative character of the system: the total energy decays monotonically as long as \(q\) is not identically zero.

\paragraph{Effect on the discrete block system}
At the semi-discrete level, the friction term modifies only the \((2,2)\) block of the spatial operator. The conservative block matrix \(\mathcal{S}\) is replaced by its damped counterpart:

\begin{equation}
\mathcal{S}_{\mathrm{fr}}
=
\begin{bmatrix}
\mathbf{0} & \mathbf{D}\\
gH\mathbf{D} & \gamma\mathbf{M}
\end{bmatrix}
\label{eq:friction_block_matrix_hw2}
\end{equation}

The block mass matrix \(\mathcal{M}\) is unchanged. Consequently, the Crank-Nicolson time discretization \eqref{eq:cn_linear_system} preserves its structure and becomes:

\begin{equation}
\left(\mathcal{M} + \frac{\Delta t}{2}\mathcal{S}_{\mathrm{fr}}\right)\underline{U}^{k+1}
=
\left(\mathcal{M} - \frac{\Delta t}{2}\mathcal{S}_{\mathrm{fr}}\right)\underline{U}^{k}
\label{eq:friction_cn_hw2}
\end{equation}

From an implementation standpoint, the friction study requires no new time integrator: it only entails the reassembly of the modified operator \eqref{eq:friction_block_matrix_hw2} before entering the time loop.

\subsection{Numerical Results and Physical Interpretation}

Two isolated numerical experiments are conducted using the Gaussian initial data \eqref{eq:gaussian_ic} on the domain \((0, \SI{1}{\meter})\) with \(H = \SI{1}{\meter}\), \(g = \SI{9.81}{\meter\per\second\squared}\) and final time \(T = \SI{0.5}{\second}\). The first experiment imposes reflecting walls with \(\gamma = 0\); the second retains periodic boundary conditions and activates linear friction with \(\gamma = 1\). In both cases, the finest mesh (\(N_e = 80\), \(p = 4\)) and a sufficiently small time step are employed to avoid temporal error contamination. The Gaussian pulse is the most suitable benchmark for these optional studies: its spatial localization makes both the reflection mechanism and the amplitude decay immediately visible in the computed profiles.

\paragraph{Reflecting walls: conservative closed-domain dynamics}

\begin{figure}[!htb]
    \centering
    \subfloat[Free-surface elevation \(\eta(x,T)\) under reflecting-wall boundary conditions \eqref{eq:reflecting_bc_hw2}. The initial rightward-propagating Gaussian pulse has reached the right wall and is reflected back into the domain. The reflected profile retains the original amplitude to a very good approximation, in agreement with the conservative energy identity \eqref{eq:energy_reflecting_conserved}]
    {\includegraphics[width=0.90\textwidth]{00_Images/HW2/P5a.1_eta_wall.png}}
    \\
    \subfloat[Discharge \(q(x,T)\) under reflecting-wall boundary conditions. The discharge goes to exactly zero at both \(x=0\) and \(x=L\), enforcing the Dirichlet constraint \eqref{eq:reflecting_bc_hw2}. This zero-discharge condition at the walls is the direct mechanism that triggers the conservative in-phase reflection of the free-surface elevation observed in the upper panel]
    {\includegraphics[width=0.90\textwidth]{00_Images/HW2/P5a.2_q_wall.png}}
    \caption{Reflecting-wall experiment at \(T = \SI{0.5}{\second}\). Physical parameters: \(L = \SI{1}{\meter}\), \(H = \SI{1}{\meter}\), \(g = \SI{9.81}{\meter\per\second\squared}\), \(\gamma = 0\). Discretization: \(N_e = 80\), \(p = 4\). The reflected wave preserves its amplitude, confirming that \(dE/dt = 0\) in a closed domain.}
    \label{fig:wall_results}
\end{figure}

The discharge profile in figure \ref{fig:wall_results}(b) shows that \(q\) is driven to exactly zero at both endpoints, directly enforcing the boundary condition \eqref{eq:reflecting_bc_hw2}. This zero-discharge constraint is the physical mechanism responsible for the in-phase reflection of \(\eta\): the wave cannot exit the domain in terms of flux, so it reverses direction while retaining its amplitude. The numerical profile in figure \ref{fig:wall_results}(a) confirms the conservation law \eqref{eq:energy_reflecting_conserved}: the reflected pulse carries the same peak height as the incident one and no energy is lost at the wall.

\paragraph{Friction: dissipative transport}

\begin{figure}[!htb]
    \centering
    \subfloat[Free-surface elevation \(\eta(x,T)\) with linear friction \(\gamma = 1\) under periodic boundary conditions. The damped solution (solid line) remains in phase with the undamped reference (dashed), confirming that the wave speed \(c = \sqrt{gH}\) is unaffected by the friction term. However, the peak amplitude is significantly reduced, visualizing the monotonic energy dissipation expressed by \eqref{eq:energy_dissipation_hw2}]
    {\includegraphics[width=0.90\textwidth]{00_Images/HW2/P5b.1_eta_friction.png}}
    \\
    \subfloat[Discharge \(q(x,T)\) with linear friction \(\gamma = 1\). The discharge profile undergoes the same amplitude attenuation as the elevation, consistently with the proportionality relation \(q_{\mathrm{ex}} = c \eta_{\mathrm{ex}}\) that holds in the undamped limit. The reduction in discharge magnitude is directly responsible for the bulk energy loss through the term \(-(\gamma/H)\int_0^L q^2 dx\)]
    {\includegraphics[width=0.90\textwidth]{00_Images/HW2/P5b.2_q_friction.png}}
    \caption{Friction experiment at \(T = \SI{0.5}{\second}\) with linear damping coefficient \(\gamma = 1\). Physical parameters and discretization as in figure \ref{fig:wall_results}. The damped wave remains in phase with the undamped reference while exhibiting a clear amplitude reduction, in agreement with the Telegrapher's equation \eqref{eq:telegrapher_hw2} and the energy inequality \eqref{eq:energy_dissipation_hw2}}
    \label{fig:friction_results}
\end{figure}

Figure \ref{fig:friction_results} shows that the linear friction term does not alter the propagation speed: the damped profile remains perfectly in phase with the undamped reference. The effect of friction is exclusively dissipative, manifesting as a progressive reduction of the wave amplitude. Both \(\eta\) and \(q\) are attenuated by the same factor, consistently with the energy law \eqref{eq:energy_dissipation_hw2}. The Telegrapher's equation \eqref{eq:telegrapher_hw2} provides the analytical explanation: the additional first-order time derivative \(\gamma\eta_t\) introduces exponential amplitude decay while leaving the wave speed \(c\) unchanged.

\paragraph{Energy analysis: validation of the analytical predictions}

\begin{figure}[!htb]
    \centering
    \includegraphics[width=0.90\textwidth]{00_Images/HW2/P5b.3_energy.png}
    \caption{Time evolution of the total discrete energy \(E(t) = \sum_i \left(\tfrac{1}{2}g\eta_i^2 + \tfrac{1}{2H}q_i^2\right)\Delta x_i\) for the Gaussian benchmark under two configurations: without friction (\(\gamma = 0\), solid line) and with linear friction (\(\gamma = 1\), dashed line). Physical parameters and discretization as in figure \ref{fig:wall_results}. The flat energy curve for \(\gamma = 0\) confirms that the combination of the SEM spatial discretization and the Crank-Nicolson time integrator preserves the discrete energy identity \eqref{eq:energy_reflecting_conserved} to machine precision. The monotonically decreasing curve for \(\gamma = 1\) is consistent with the dissipation law \eqref{eq:energy_dissipation_hw2}, providing quantitative numerical validation of the analytical derivation}
    \label{fig:energy_evolution}
\end{figure}

Figure \ref{fig:energy_evolution} provides the conclusive numerical validation of the analytical derivations carried out in sections \ref{ssc:251} and \ref{ssc:252}. In the frictionless case (\(\gamma = 0\)), the total discrete energy remains constant throughout the simulation: the flat line confirms that the Crank-Nicolson scheme, combined with the SEM spatial discretization, preserves the conservative energy identity \eqref{eq:energy_reflecting_conserved} to machine precision. This is a non-trivial property: it reflects the symplectic-like character of the Crank-Nicolson integrator when applied to a skew-symmetric spatial operator. When friction is activated (\(\gamma = 1\)), the energy curve decreases monotonically from its initial value, in full agreement with the theoretical inequality \eqref{eq:energy_dissipation_hw2}. The rate of decay is governed by the bulk dissipation term \(-(\gamma/H)\int_0^L q^2  dx\), which is largest when the discharge is at its maximum and tapers off as the wave amplitude diminishes over time.